\documentclass[cs4size,a4paper]{ctexart}
\usepackage{xeCJK}

%章节总标题加点
\usepackage[subfigure]{tocloft}	
\renewcommand{\cftsecleader}{\cftdotfill{\cftdotsep}}  

%==================== 数学符号公式 ============
\usepackage{amsmath}                 % AMS LaTeX宏包
%\usepackage{amssymb}                % 用来排版漂亮的数学公式
%\usepackage{amsbsy}
\usepackage[style=1]{mdframed}
\usepackage{amsthm}
\usepackage{amsfonts}
\usepackage{mathrsfs}                % 英文花体字体
\usepackage{bm}                      % 数学公式中的黑斜体
\usepackage{bbding,manfnt}           % 一些图标，如 \dbend
\usepackage{lettrine}                % 首字下沉，命令\lettrine
\def\attention{\lettrine[lines=2,lraise=0,nindent=0em]{\large\textdbend\hspace{1mm}}{}}
\usepackage{longtable}
\usepackage[toc,page]{appendix}
\usepackage{geometry}                % 页边距调整
\geometry{top=3.0cm,bottom=2.7cm,left=2.5cm,right=2.5cm}
%\usepackage{caption2}               % 浮动图形和表格标题样式

%====================按章编号==========================
\numberwithin{equation}{section}
\numberwithin{table}{section}
\numberwithin{figure}{section}

%================= 基本格式预置 ===========================
\linespread{1.25} \zihao{-4}  % 行距、字号
\setlength{\parskip}{0pt}  % 段距

% 页眉页脚
\usepackage{fancyhdr}
\pagestyle{fancy}
\fancyhf{}  
\fancyhead[C]{\zihao{5} \songti 西安石油大学本科毕业设计（论文）}
\fancyfoot[C]{~\zihao{5} —\hspace{0.5em} \thepage \hspace{0.5em}—~}

\renewcommand{\headrulewidth}{0.65pt} 

% 标题格式
\CTEXsetup[format={\centering\zihao{-3}\bfseries},beforeskip={1em},afterskip={2em}]{section}
\CTEXsetup[format={\bfseries\zihao{4}},beforeskip={0.5em},afterskip={0.5em}]{subsection}
\CTEXsetup[format={\bfseries\zihao{-4}},beforeskip={0.5em},afterskip={0.5em}]{subsubsection}

% 上标引用格式
\newcommand{\upcite}[1]{\textsuperscript{\textsuperscript{\cite{#1}}}}  
\renewcommand{\cftsecleader}{\cftdotfill{\cftdotsep}}
\usepackage[font=small,labelfont=bf, textfont=bf]{caption}

% 目录不加粗
%\usepackage{tocstyle}
%\usepackage{microtype}
%\settocfeature[toc][1]{entryhook}{\normalfont}


%================== 图形支持宏包 =========================
\usepackage{subfigure}
\usepackage{graphicx}                % 嵌入png图像
\usepackage{color,xcolor}            % 支持彩色文本、底色、文本框等
\usepackage{pifont}                  % 支持圈数字 %\ding{number}
\usepackage{hyperref}                % 交叉引用
\usepackage{caption}
\captionsetup{figurewithin=section}

%==================== 源码和流程图 ====================
\usepackage{listings}                % 粘贴源代码
\lstset{ 
	language=PHP, 
	tabsize=4, 
	basicstyle=\normalsize\linespread{1.1}\footnotesize\ttfamily,
	showstringspaces=false,
	frame=top, frame=bottom,
	commentstyle=\color{green}, 
	keywordstyle=\color{blue}, 
	breaklines=true,
	numbers=left, numberstyle=\tiny,
}
\usepackage{tikz}                    
\usepackage{tikz-3dplot}
\usetikzlibrary{shapes,arrows,positioning}

%==================== 字体设置 ====================
% Mac 字体设置
%\setCJKmainfont[BoldFont=宋体-简 粗体,ItalicFont=宋体-简 常规体]{宋体-简 常规体}
%\setCJKsansfont[BoldFont=宋体-简 常规体]{宋体-简 常规体}
%\setCJKmonofont{宋体-简 常规体}

\setmainfont{Times New Roman} 
\setCJKmainfont[BoldFont=华文中宋]{宋体}

%==================== 自定义包 =======================

%===================  定理类环境定义 ===================
\newtheorem{example}{例}              % 整体编号
\newtheorem{algorithm}{算法}
\newtheorem{theorem}{定理}            % 按 section 编号
\newtheorem{definition}{定义}
\newtheorem{axiom}{公理}
\newtheorem{property}{性质}
\newtheorem{proposition}{命题}
\newtheorem{lemma}{引理}
\newtheorem{corollary}{推论}
\newtheorem{remark}{注解}
\newtheorem{condition}{条件}
\newtheorem{conclusion}{结论}
\newtheorem{assumption}{假设}
\usepackage{enumitem}


%===================   正文开始   ===================

\begin{document}
	
	%==================重定义 ===================
	\renewcommand{\contentsname}{\hspace*{\fill}目\quad 录\hspace*{\fill}\\}   
	\renewcommand{\abstractname}{摘要} 
	\renewcommand{\refname}{参考文献}     
	\renewcommand{\indexname}{索引}
	\renewcommand\thefigure{\thesection-\arabic{figure}}
	\renewcommand{\figurename}{图}
	\renewcommand\thetable{\thesection-\arabic{table}}
	\renewcommand{\tablename}{表}
	\renewcommand{\appendixname}{附录}
	\renewcommand{\proofname}{证明}
	\renewcommand{\algorithm}{算法} 
	%\renewcommand{\labelenumi}{(\arabic{enumi})}

	% 封皮和前言
	%===============  封面  =================
\smallskip
\vspace*{2.0cm}
\begin{center}
\begin{figure}[!th]
\centering
\includegraphics[width=0.3\linewidth]{figure/SchoolName}
\end{figure}

\vspace*{1.0cm}{\zihao{-1} \textbf{\songti{} 本科毕业设计（论文）}} \\
\vspace*{9.0cm} \songti{}
% 创建表内换行
\newcommand{\tabincell}[2]{\begin{tabular}{@{}#1@{}}#2\end{tabular}}
\begin{tabular}{lc}
  \zihao{4} \textbf{题\hspace{\fill}目：} & 
  \underline{\makebox[7cm][c]{\tabincell{c}{\zihao{4}海量数据的异常项检测\\ 
  \zihao{4}及分类方法研究}}} \\ 
\\
  \zihao{4}\textbf{学生姓名：} & 
  \underline{\makebox[7cm][c]{\zihao{4}刘阳}} \\ 
\\
  \zihao{4}\textbf{院\hspace{\fill}（系）：} & 
  \underline{\makebox[7cm][c]{\zihao{4}计算机学院}} \\ 
\\
  \zihao{4}\textbf{专业班级：} & 
  \underline{\makebox[7cm][c]{\zihao{4}软件1302}} \\ 
\\
  \zihao{4}\textbf{指导教师：} & 
  \underline{\makebox[7cm][c]{\zihao{4}张留美}} \\ 
\\
  \zihao{4}\textbf{完成时间：} & 
  \underline{\makebox[7cm][c]{\zihao{4}2017年6月}} \\ 
\\
\end{tabular} 
\end{center}
\thispagestyle{empty}
\clearpage

	\pagestyle{plain} 		%不带页眉横线
	\thispagestyle{empty} 
\section*{\zihao{3} \centering 海量数据的异常项检测及分类方法研究}
\vspace{0.5cm}

\noindent 
\textbf{摘\hspace{1em}要}：随着互联网的发展，人们正处于一个信息爆炸的时代。相比于过去的信息匮乏，面对现阶段海量的信息数据，人们淹没在数据中难以快速的制定合适的决策，使用简单的数据处理办法已经不能适用于现如今的需求。

 针对现状，本文采用不同数据处理方法来进行研究并对实现对给定的数据集进行数据检测分析，寻找能够在相同的硬件环境下提高海量数据处理效率的算法。提出了基于决策树的数据处理方法，即对数据进行特征提取、原始样本数据的预处理、构造决策树模型、进行数据挖掘预测等方面进行了深入讨论，并详细阐述了运用算子进行数据集处理的具体实现过程。利用rapidminer对数据处理过程进行可视化解析。通过对数据处理分析，所得数据具有较强的可读性、较好的可预测性。

设计并完成了数据在Web端的录入及储存在mysql数据库中，采用可视化工具rapidminer对mysql数据库中的数据进行读入，然后进行对读入数据进行预处理，即完成异常数据的检测及排除，在完成预处理后进一步进行所需属性的筛选，接着设置Label规则最后将处理后的数据传送给决策树模型，产生决策树可视化图像并最终根据决策树传出数据作出预测。

采用本文的数据处理方法可以大幅度提高数据使用效率，直观立体的反应海量数据对未来预测的指导性作用，明显优于传统数据处理方法。

\vspace{0.5cm}
\noindent
\textbf{关键词：}海量数据 ；数据处理；推荐算法；属性聚类；决策树
%\addcontentsline{toc}{section}{摘要}%摘要是否包括进目录中

\clearpage

\thispagestyle{empty} 
\section*{\songti\zihao{3} \centering \textbf{Research on the Detection and Classification of Exception in Massive Data
}}
   %用了Times New Roman字体来美化观感

\noindent 
\textbf{Abstract: }With the development of the Internet, people are in an era of information explosion. Compared to the lack of information in the past, in the face of massive information at this stage, it is difficult to develop appropriate decisions in the data, and the use of simple data processing can not be applied to today's needs.


In this paper, different data processing methods are used to study and realize the data detection and analysis of a given data set, and find an algorithm which can improve the efficiency of massive data processing in the same hardware environment. The data processing method based on decision tree is proposed, which is the feature extraction of the data, the preprocessing of the original sample data, the construction of the decision tree model, the data mining prediction and so on, and elaborates the use of the operator to carry out the data set Processing the specific implementation process. Using RapidMiner to visualize the data processing process. Through the data processing analysis, the obtained data has a strong readability, better predictability.

Design and complete the data in the Web side of the input and stored in the mysql database, the use of visual tools RapidMiner mysql database to read the data, and then read the data pre-processing, that is, the completion of abnormal data detection and exclusion, in the After the preprocessing, the filtering of the required attributes is carried out, and then the Label rule is set. Finally, the processed data is sent to the decision tree model, the decision tree visualization image is generated and finally the data is predicted based on the decision tree.

The data processing method of this paper can greatly improve the efficiency of data utilization, and the guidance function of intuitive stereo response to future forecasting is obviously superior to traditional data processing method.


\noindent 
\textbf{Keywords: }Massive Data; Data Processing; Recommendation Algorithm; Attribute Clustering;Decision tree
%\addcontentsline{toc}{section}{Abstract}%英文摘要是否包括进入目录

	
	
	%=================== 目录 ===================
	% 目录页加页眉
	\fancypagestyle{plain}{  
		\fancyhead[C]{\zihao{5} \songti 目\hspace{1em}录}
	}
	\pagenumbering{Roman}
	
	\tableofcontents 
	
	
	%============== 论文正文   =================
	%\setcounter{tocdepth}{4}
	%\setcounter{secnumdepth}{4}
	\pagestyle{fancy}
	\pagenumbering{arabic}
\section{绪论}

\subsection{研究背景}

随着互联网的高速发展，人们生活在数字化的信息时代。相比以前较难获取信息，当下的人们可以十分便捷地接触到海量数据，以至于人们淹没在数据中难以快速的制定合适的决策。在校园中，学院的日常教务工作同样面临着这个难题，迫切地需要将已有的业务逻辑线上化。

其中，西安石油大学的毕业设计选题、中期检查、答辩等主要教学工作和学生工作，基本采用人工或半信息化的方式进行管理。针对这些陈旧的毕业管理工作，利用了新型网络技术，设计并实现毕业管理云平台系统，解决传统管理方式的效率低、工作量大、易出错等关键问题。

XSYU-GMS 系统是一个用于高校毕业生毕业流程线上管理的教务系统，由笔者独立开发，并于 2016 年起在西安石油大学计算机学院投入运行，承担了学院毕业设计选题工作。经过两年学院内测试后，该系统已基本完善，并且工作状态良好。本项目将在 GMS 系统的基础上进行二次开发，对接之前实现的毕业设计选题功能，实现毕业设计答辩系统。

\subsection{研究需求与意义}

本项目将使用新技术实现毕业设计答辩系统，解决在线下工作的不便以及存在的各种问题，同时新增大量特性以降低使用者的工作量。无论对于学生、教师还是管理人员，这套新系统都将会提供比以前更稳定易用的工作方式。

该项目使用了目前较为流行的 React 框架和 Node.js 后端技术，这样能使学生熟悉掌握目前比较新的 Web 开发技术。通过答辩系统项目的实践开发，学生可以混合使用在本科所学的多门学科，在实践中合理利用理论知识，提高学生在程序开发上的的动手能力。

本线上管理系统采用目前最新的 Web 开发技术搭建而成，具有可扩展、高可用性的特点，可以在后期根据需求开发新的模块。注重用户体验，完全按照学院规定的教学流程开发，方便教师学生使用。在之后可为西安石油大学计算机学院的教学管理工作定制开发，成为学院内与教务系统平行的一款教学辅助系统。

\subsection{研究现状及其前景}

几乎所有的高校都需要面对毕业生毕业设计、答辩的相关工作，基本没有学校再采用人工的方式进行管理。然而目前市面上并没有一款专门为此目的而开发的商业系统可供直接使用，大多数院校通过软件公司开发，或自己开发，这难免会造成功能不满足要求，或者开发质量低等问题。

部分学校所选用的教务系统或许有关于毕业设计的子模块，不过功能简单，还不足以满足学院教学工作的需要。目前对于毕业设计答辩流程的线上管理，将主要还是由各高校根据自己的相关需求，独立开发或外包开发。

类似的项目有阿尔巴尼亚高中设计的一个复杂的管理信息系统。该应用程序为其用户操作提供了一个合适的接口，以便简化和减少管理各个学校活动有关的信息和过程所需的时间。随着教育管理体制的深化和信息技术的进步，教育管理呈现出系统管理与教学管理脱节的趋势。该开发者认为，其根本原因在于教育管理体制整体建设的缺失。而山东师范大学信息科学与工程学院独立自主开发的“学生信息管理系统”软件，用于管理学生的学术数据。类似的，Norasiah et. al. 设计了一个学生档案管理系统 (SRMS) 信息系统 ，提供了一个直观的界面来管理学生数据。它能被很多大学或学院使用，以有效地记录学生的数据。它使用了序列图和状态图来详细分析数据的主要情况，然后给出系统的动态模型。并在教学管理系统面向对象建模的基础上，提出了一种设计过程。

尽管上述研究显示了一些让传统学生管理行业或行政管理所受益的项目，但仍不足以满足这种商业逻辑下的现代化要求。因此，我们需要开发一个更可靠的、更高性能的信息系统。


\subsection{传统 Web 开发技术简介}

在传统 Web 开发项目中，CGI（Common Gateway Interface）是能让 Web 服务器和 CGI 脚本正确处理客户的请求的协议。CGI 可以通过 C、PHP、Python、JSP 等多种编程语言实现。

如图 1-1 所示。其中 Web 服务器负责管理和维护来自用户的连接请求、基于 HTTP 协议的数据传输，以及网络交互等。CGI 程序则提供对具体的业务逻辑的实现。Web 服务器的功能主要是将客户端发来的 HTTP 请求转换成相应的 CGI 程序请求，然后执行目标程序，接着将 CGI 程序的输出作为结果，发送给客户端实现信息回复。

\begin{figure}
	\centering
	\includegraphics[width=0.7\linewidth]{figure/1-1}
	\caption{传统 Web 项目实现方式逻辑图}
\end{figure}

传统的 Web 开发，那都是由前端发送一个 HTTP 请求，一般包括 GET 请求和 POST 请求两种。后端处理该请求，并整理出 HTML 页面再返回。AJAX 也差不多，同样要经过发送请求、接收数据、写入页面的过程。应用逻辑都在后端，无数的请求往往成为开发的中心，但对于业务而言其实却是多余的。

除了数据传输时需要的封装请求和解析请求操作较为繁琐。我们还可以发现，页面的 UI 在后端产生，但又需要配合前端 CSS 进行维护，也可以在前端进行页面的二次修改，页面的显示逻辑分布在了前后端两侧，为项目维护造成了麻烦。

类似的情况发生在业务逻辑上，传统 Web 项目主要业务逻辑在后端完成。然而随着前端 JavaScript 的快速发展，很多简单逻辑被放在前端中执行。这虽然提升了项目开发效率，但同样带来了业务逻辑前后端分离的问题，在中大型项目中很容易会带来混乱，使项目不可维护。

RESTful（具象状态传输）架构是 Web 开发进一步发展后的成果，如图 1-2 所示。这是一种架构风格，它定义了一组基于HTTP的约束和属性。符合REST架构风格的Web服务提供了互联网上计算机系统之间的相互可操作性。符合REST的Web服务允许请求系统通过使用统一和预定义的一组无状态操作，来访问和操纵Web资源的文本表示（一般为 JSON 形式）。其他类型的Web服务有SOAP Web服务等。前端在获取静态的 HTML 页面之后，所有的业务逻辑和数据的获取操作均使用后端提供的 API 实现。在后端 API 保持不变的情况下，大部分业务逻辑在前端实现。这样使得项目开发更加规则，易于维护。但事实上，这种方式仍有少部分业务逻辑需要在后端实现。

\begin{figure}
	\centering
	\includegraphics[width=0.7\linewidth]{figure/1-2}
	\caption{RESTful 架构开发方式}
\end{figure}

而本次答辩系统开发采用了更彻底地逻辑分离，在前后端分离的基础上抛弃了后端 API，几乎是直接暴露了后端数据库操作权限。消除了传统 Web 开发模式当中的请求封装、请求解析、拼装页面等一系列繁琐且多余的操作，集中在清晰的业务逻辑中，这有益于项目的快速开发和维护。

      
	\pagenumbering{arabic}
\section{异常检查}
支持向量机（Support Vector Machies）是由Vapnik等人于1995年提出来的。之后随着统计理论的发展，支持向量机也逐渐受到了各领域研究者的关注，在很短的时间就得到很广泛的应用。支持向量机是建立在统计学习理论的VC维理论和结构风险最小化原理基础上的，利用有限的样本所提供的信息对模型的复杂性和学习能力两者进行了寻求最佳的折衷，以获得最好的泛化能力。SVM的基本思想是把训练数据非线性的映射到一个更高维的特征空间（Hilbert空间）中，在这个高维的特征空间中寻找到一个超平面使得正例和反例两者间的隔离边缘被最大化。SVM的出现有效的解决了传统的神经网络结果选择问题、局部极小值、过拟合等问题。并且在小样本、非线性、数据高维等机器学习问题中表现出很多令人注目的性质，被广泛地应用在模式识别，数据挖掘等领域(张学工 2000；崔伟东2001)。支持向量机可以用于分类和回归问题，本章着重介绍分类相关的知识。\\

\subsection{离群点检测}


就餐饮企业而言，经常会碰到这样的问题：

\begin{itemize}[leftmargin=1cm]

\item 如何根据客户的消费记录检测是否为异常刷卡消费？

\item 如何检测是否有异常订单？

\end{itemize}

这一类异常问题可以通过离群点检测解决。

离群点检测是数据挖掘中重要的一部分，它的任务是发现与大部分其他对象显著不同的对象。大部分数据挖掘方法都将这种差异信息视为噪声而丢弃，然而在一些应用中，罕见的数据可能蕴含着更大的研究价值。

在数据的散布图中，如图2-1离群点远离其它数据点。因为离群点的属性值明显偏离期望的或常见的属性值，所以离群点检测也称偏差检测。

\begin{figure}[thbp!]
\centering
\includegraphics[width=0.4\linewidth]{figure/2-1}
\includegraphics[width=0.4\linewidth]{figure/2-2}
\caption{\ 离群点检测示意图}
\label{fig:2-2}
\end{figure}

%latex自动控制页面布局，一般你不需要管页面
离群点检测已经被广泛应用于电信和信用卡的诈骗检测、贷款审批、电子商务中、网络入侵、天气预报等领域，如可以利用离群点检测分析运动员的统计数据，以发现异常的运动员。

\begin{itemize}[leftmargin=1cm]

\item 离群点的成因 离群点的主要成因有：数据来源于不同的类、自然变异、数据测量和收集误差。

\item 离群点的类型  对离群点的大致分类见表2-1
\end{itemize}

\begin{table}[thbp]\footnotesize
\caption{离群点的大致分类}
\begin{center}
\begin{tabular}{c|cl}
\hline  分类标准& 分类名称 & 分类描述  \\ 
\hline 从数据范围 & 全局离群点和局部离群点 & 从整体来看，某些对象没有离群特征， \\ 
& & 但是从局部来看，却显示了一定的离群性。\\ 
& & 如图2-2：C是全局离群点，D是局部离群点。 \\ 
\hline 从数据类型 & 数值型离群点和分类型离群点 & 这是以数据集的属性类型进行划分的。  \\ 
\hline 属性的个数 & 一维离群点和多维离群点 & 一个对象可能有一个或多个属性。 \\ 
\hline
\end{tabular} 
\end{center}
\label{tb:filter}
\end{table}

\begin{figure}[thbp!]
\centering
\includegraphics[width=0.4\linewidth]{figure/2-3}
\caption{\ 全局离群点和局部离群点}
\label{fig:2-3}
\end{figure}

%\end{enumerate}


\subsubsection{离群点检测方法}

常用离群点检测方法见表\ref{tb:filter1}

\begin{table}[thbp]%表格你自己再调整，很简单，就是再插入行，然后不加hline，一列不放数据
\footnotesize
\caption{常用离群点检测方法}
\begin{center}
\begin{tabular}{|l|l|l|}
\hline  离群点检测方法& 方法描述 & 方法评估  \\ 
\hline 基于统计 & 大部分的基于统计的离群点检& 基于统计模型的离群点检测方法 \\  
&测方法是构建一个概率分布模型， &的前提是必须知道数据集服从什么分布； \\
&并计算对象符合该模型的概率，&对于高维数据，检验效果可能很差。\\
&把具有低概率的对象视为离群点。& \\
\hline 基于邻近度 & 通常可以在数据对象之间定义邻近性度量， & 简单，二维或三维的数据可以做散点图观察；  \\ 
& 把远离大部分点的对象视为离群点。&大数据集不适用；对参数选择敏感；具有全局\\
& & 阈值，不能处理具有不同密度区域的数据集。\\
\hline 基于密度 & 考虑数据集可能存在不同密度区域这一事实， & 给出了对象是离群点的定量度量， \\ 
&从基于密度的观点分析，离群点是在低密度区域中的对象。&并且即使数据具有不同的区域也能够很好的处理；\\
&一个对象的离群点得分是该对象周围密度的逆。&大数据集不适用；参数选择是困难的。\\
\hline 基于聚类 & 一种是利用聚类检测离群点的方法是丢弃远离其他簇的小簇；& 基于聚类技术来发现离群点可能是高度有效的； \\ 
&另一种更系统的方法，首先聚类所有对象，&聚类算法产生的簇的质量对该算法产生\\
&然后评估对象属于簇的程度（离群点得分）。 &的离群点的质量影响非常大。\\
\hline

\end{tabular} 
\end{center}
\label{tb:filter1}%label 是唯一的
\end{table}

基于统计模型的离群点检测方法需要满足统计学原理，如果分布已知，则检验可能非常有效。基于邻近度的离群点检测方法比统计学方法更一般、更容易使用，因为确定数据集有意义的邻近度量比确定它的统计分布更容易。基于密度的离群点检测与基于邻近度的离群点检测密切相关，因为密度常用邻近度定义：一种是定义密度为到K个最邻近的平均距离的倒数，如果该距离小，则密度高；另一种是使用DBSCAN聚类算法，一个对象周围的密度等于该对象指定距离d内对象的个数。

本节重点介绍基于统计模型和聚类的离群点检测方法。
\subsubsection{基于模型的离群点检测方法}
通过估计概率分布的参数来建立一个数据模型，如果一个数据对象不能很好地跟该模型拟合，即如果它很可能不服从该分布，则它是一个离群点。

%大段落不能用itemize
%\begin{itemize}[leftmargin=1cm]

(1) 一元正态分布中的的离群点检测

正态分布是统计学中最常用的分布之一。

若随机变量$x$的密度函数$\varphi(x)=\frac{1}{\sqrt{2\pi}}e\frac{-(x-\mu)^2}{2\sigma^2}(x\in R)$，则称$x$服从正态分布，简称$x$服从正态分布$N(\mu,\sigma)$，其中参数$\mu$和$\sigma$分别为均值和标准差。

图2-3显示$N(0,1)$的密度函数：
\begin{figure}[thbp!]
\centering
\includegraphics[width=0.4\linewidth]{figure/2-6}
\caption{$N(0,1)$的概率密度函数}
\label{fig:2-6}
\end{figure}

$N(0,1)$的数据对象出现在该分布的两边尾部的机会很小，因此可以用它作为检测数据对象是否是离群点的基础。数据对象落在三倍标准差中心区域之外的概率仅有0.0027。

(2)混合模型的离群点检测

这里首先介绍下混合模型。混合是一种特殊的统计模型，它使用若干统计分布对数据建模。每一个分布对应一个簇，而每个分布的参数提供对应簇的描述，通常用中心和发散描述。

混合模型将数据看作从不同的概率分布得到的观测值的集合。概率分布可以是任何分布，但是通常是多元正态的，因为这种类型的分布不难理解，容易从数学上进行处理，并且已经证明在许多情况下都能产生好的结果。这种类型的分布可以对椭圆簇建模。

总的讲，混合模型数据产生过程为：给定几个类型相同但参数不同的分布，随机地选取一个分布并由它产生一个对象。重复该过程次，其中是对象的个数。

具体地讲，假定有$K$个分布和$m$个对象$\chi=\{x_1,x_2,...,x_m\}$。设第$ j$个分布的参数为$\alpha_j$，并$A$设是所有参数的集合，即$A=\{\alpha_1,\alpha_2,...,\alpha_k\}$。则$p(x_i|a_j)$是第$i$个对象来自第$j$个分布的概率。选取第$j$个分布产生一个对象的概率由权值$w_j\{1\le j\le K\}$给定，其中权值（概率）受限于其和为1的约束，即$\sum_{j=1}^k w_j=1$。于是，对象$x$的概率由以下公式给出：
\begin{equation}p(x|A)=\sum_{j=1}^k w_jP_j(x|\theta_j)\end{equation}
如果对象以独立的方式产生，则整个对象集的概率是每个个体对象$x_i$的概率的乘积，公式如下：
\begin{equation}p(\chi|\alpha)=\prod_{i=1}^m P(x_i|\alpha)=\prod_{i=1}^m\sum_{j=1}^k w_jP_j(x|\alpha_j)\end{equation}
对于混合模型，每个分布描述一个不同的组，即一个不同的簇。通过使用统计方法，可以由数据估计这些分布的参数，从而描述这些分布（簇）。也可以识别哪个对象属于哪个簇。然而，混合模型只是给出具体对象属于特定簇的概率。

聚类时，混合模型方法假定数据来自混合概率分布，并且每个簇可以用这些分布之一识别。同样，对于离群点检测，数据用两个分布的混合模型建模，一个分布为正常数据，而另一个为离群点。

聚类和离群点检测的目标都是估计分布的参数，以最大化数据的总似然。

这里提供一种离群点检测常用的简单的方法：先将所有数据对象放入正常数据集，这时离群点集为空集；再用一个迭代过程将数据对象从正常数据集转移到离群点集，只要该转移能提高数据的总似然。

具体操作如下：

假设数据集U包含来自两个概率分布的数据对象：M是大多数（正常）数据对象的分布，而N是离群点对象的分布。数据的总概率分布可以记作：

 $U(x)=(1-\lambda)M(x)+\lambda N(x)$其中，$x$是一个数据对象；$\lambda\in[0,1]$，给出离群点的期望比例。分布$M$由数据估计得到，而分布$N$通常取均匀分布。设$M_t$和$N_t$分别为时刻$t$正常数据和离群点对象的集合。初始$t=0$，$M_0=D$，而$N_0=\emptyset$。

根据混合模型中公式$p(x|A)=\sum_{j=1}^k w_jP_j(x|\theta_j)$推导，在整个数据集的似然和对数似然可分别由下面两式给出：
\begin{equation} 
L_t(U)=\prod_{x_i\in U } P_U(x_i)=\bigg ((1-\lambda)^{|M_t|} \prod_{x_i\in M_i} P_{M_i}(x_i) \bigg)\bigg(\lambda^{|N_t|} \prod_{x_i\in N_i}P_{N_i}(x_i)\bigg) 
\end{equation}
\begin{equation}
\ln L_t(U)=|M_t|\ln (1-\lambda)+\sum_{x_i \in M_i }\ln P_{M_i}(x_i)+|N_t|\ln \lambda +\sum_{x_i \in N_i }\ln P_{N_i}(x_i)
\end{equation}
其中$P_D$、$P_{M_t}$、$P_{N_t}$分别是$D$、$M_t$、$N_t$的概率分布函数。

因为正常数据对象的数量比离群点对象的数量大的很多，因此当一个数据对象移动到离群点集后，正常数据对象的分布变化不大。在这种情况下，每个正常数据对象的正常数据对象的总似然的贡献保持不变。此外，如果假定离群点服从均匀分布，则移动到离群点集的每一个数据对象对离群点的似然贡献一个固定的量。这样，当一个数据对象移动到离群点集时，数据总似然的改变粗略地等于该数据对象在均匀分布下的概率（用$1-\lambda$加权）减去该数据对象在正常数据点的分布下的概率（用加权）。从而，离群点由这样一些数据对象组成，这样数据对象在均匀分布下的概率比正常数据对象分布下的概率高。

在某些情况下是很难建立模型的。如：因为数据的统计分布未知或没有训练数据可用。在这种情况下，可以考虑另外其他不需要建立模型的检测方法。

\subsubsection{基于聚类的离群点检测方法}
聚类分析用于发现局部强相关的对象组，而异常检测用来发现不与其他对象强相关的对象。因此聚类分析非常自然地可以用于离群点检测。本节主要介绍两种基于聚类的离群点检测方法。



(1)丢弃远离其他簇的小簇

一种利用聚类检测离群点的方法是丢弃远离其他簇的小簇。通常，该过程可以简化为丢弃小于某个最小阈值的所有簇。

这个方法可以和其他任何聚类技术一起使用，但是需要最小簇大小和小簇与其他簇之间距离的阈值。而且这种方案对簇个数的选择高度敏感，使用这个方案很难将离群点得分附加到对象上。

图2-4中，聚类簇数K=2，可以直观地看出其中一个包含5个对象的小簇远离大部分对象，可以视为离群点。

\begin{figure}[thbp!]
\centering
\includegraphics[width=0.4\linewidth]{figure/2-7}
\caption{\ K-Means算法的聚类图}
\label{fig:2-7}
\end{figure}

(2)基于原型的聚类

另一种更系统的方法，首先聚类所有对象，然后评估对象属于簇的程度（离群点得分）。在这种方法中，可以用对象到它的簇中心的距离来度量属于簇的程度。特别地，如果删除一个对象导致该目标的显著改进，则可将该对象视为离群点。例如，在K均值算法中，删除远离其相关簇中心的对象能够显著地改进该簇的误差平方和（SSE）。

对于基于原型的聚类，评估对象属于簇的程度（离群点得分）主要有两种方法：一是度量对象到簇原型的距离，并用它作为该对象的离群点得分；二是考虑到簇具有不同的密度，可以度量簇到原型的相对距离，相对距离是点到质心的距离与簇中所有点到质心的距离的中位数之比。

如图2-5，如果选择聚类簇数K=3，则对象A、B、C应分别属于距离它们最近的簇，但相对于簇内的其他对象，这三个点又分别远离各自的簇，所以有理由怀疑对象A、B、C是离群点。


\begin{figure}[thbp!]
\centering
\includegraphics[width=0.4\linewidth]{figure/2-8}
\caption{\ 基于距离的离群点检测}
\label{fig:2-8}
\end{figure}


诊断步骤如下：


(1)进行聚类。

选择聚类算法（如K-Means算法），将样本集聚为K簇，并找到各簇的质心。

(2)计算各对象到它的最近质心的距离。

(3)计算各对象到它的最近质心的相对距离。

(4)与给定的阈值作比较。

如果某对象距离大于该阈值，就认为该对象是离群点。



基于聚类的离群点检测的改进：



(1)离群点对初始聚类的影响：通过聚类检测离群点时，离群点会影响聚类结果。为了处理该问题，可以使用如下方法：对象聚类，删除离群点，对象再次聚类（这个不能保证产生最优结果）。

(2)还有一种更复杂的方法：取一组不能很好的拟合任何簇的特殊对象，这组对象代表潜在的离群点。随着聚类过程的进展，簇在变化。不再强属于任何簇的对象被添加到潜在的离群点集合；而当前在该集合中的对象被测试，如果它现在强属于一个簇，就可以将它从潜在的离群点集合中移除。聚类过程结束时还留在该集合中的点被分类为离群点（这种方法也不能保证产生最优解，甚至不比前面的简单算法好，在使用相对距离计算离群点得分时，这个问题特别严重）。

对象是否被认为是离群点可能依赖于簇的个数（如k很大时的噪声簇）。该问题也没有简单的答案。一种策略是对于不同的簇个数重复该分析。另一种方法是找出大量小簇，其想法是：

\begin{itemize}[leftmargin=1cm]

\item 较小的簇倾向于更加凝聚；

\item 如果存在大量小簇时一个对象是离群点，则它多半是一个真正的离群点。

不利的一面是一组离群点可能形成小簇从而逃避检测。
\end{itemize}



“Detect Outlier(Distances)”基于距离的离群点检测，参数设置中可设定要检测的离群点的个数，如图2-6。

\begin{figure}[thbp!]
\centering
\includegraphics[width=0.4\linewidth]{figure/2-9}
\caption{\  RapidMiner自带的离群点检测流程}
\label{fig:2-9}
\end{figure}


第三方离群点检测插件带有功能更强的离群点检测功能，例如“One-Class LIBSVM Anomaly Score”为半监督的离群点检测操作符。

 离群点检测实例

下面，我们自己生成一个数据，来看看离群点检测的功能。

第一步：生成随机数据

调用“Generate Data”生成数据操作符，能帮助我们自动创建一些测试数据，创建参数设置图2-7。

\begin{figure}[thbp!]
\centering
\includegraphics[width=0.4\linewidth]{figure/2-10}
\caption{\ 生成随机数据参数设置}
\label{fig:2-10}
\end{figure}


调用“Map”映射操作符，设置参数如图2-8，将所有的数据类型都转换为normal类型。

\begin{figure}[thbp!]
\centering
\includegraphics[width=0.4\linewidth]{figure/2-11}
\caption{\  映射操作符参数设置}
\label{fig:2-11}
\end{figure}


再次调用“Generate Data”生成数据操作符，参数设置如图2-9，添加离群点

\begin{figure}[thbp!]
\centering
\includegraphics[width=0.4\linewidth]{figure/2-12}
\caption{\  添加离群点参数设置}
\label{fig:2-12}
\end{figure}


同样，添加Map操作符，参数设置如图2-10

\begin{figure}[thbp!]
\centering
\includegraphics[width=0.4\linewidth]{figure/2-13}
\caption{\ Map映射离群点设置}
\label{fig:2-13}
\end{figure}


最后添加“Append”操作符将两个数据集合并，全部操作流程图如图2-11，数据散点图如图2-12.

\begin{figure}[thbp!]
\centering
\includegraphics[width=0.4\linewidth]{figure/2-14}
\caption{\ 1.操作流程图}
\label{fig:2-14}
\end{figure}


\begin{figure}[thbp!]
\centering
\includegraphics[width=0.4\linewidth]{figure/2-15}
\caption{\ 2.数据散点图}
\label{fig:2-15}
\end{figure}

第二步：离群点检测操作符应用

“k-NN Global Anomaly Score”k-NN全局离群点检测操作符，检测结果如图2-13

\begin{figure}[thbp!]
\centering
\includegraphics[width=0.4\linewidth]{figure/2-16}
\caption{\ 全局离群点检测气泡图}
\label{fig:2-16}
\end{figure}

“Local Outlier Factor”基于本地的离群点检测操作符，操作流程如图2-14，检测结果如图2-15

\begin{figure}[thbp!]
\centering
\includegraphics[width=0.4\linewidth]{figure/2-17}
\caption{\ 离群点检测操作流程}
\label{fig:2-17}
\end{figure}

\begin{figure}[thbp!]
\centering
\includegraphics[width=0.4\linewidth]{figure/2-18}
\caption{\ 离群点检测操作流程}
\label{fig:2-18}
\end{figure}


	\pagenumbering{arabic}
\section{决策树算法实现流程}
\subsection{什么是决策树算法}
决策树是一树状结构，它的每一个叶节点对应着一个分类，非叶节点对应着在某个属性上的划分，根据样本在该属性上的不同取值将其划分成若干个子集。对于非纯的叶节点，多数类的标号给出到达这个节点的样本所属的类。构造决策树的核心问题是在每一步如何选择适当的属性对样本做拆分。对一个分类问题，从已知类标记的训练样本中学习并构造出决策树是一个自上而下，分而治之的过程。

本节下面的内容将开始讲解决策树的具体内容。首先将从一个最经典的决策树分类算法开始说明决策树。

ID3算法基于信息熵来选择最佳测试属性。它选择当前样本集中具有最大信息增益值的属性作为测试属性；样本集的划分则依据测试属性的取值进行，测试属性有多少不同取值就将样本集划分为多少子样本集，同时决策树上相应于该样本集的节点长出新的叶子节点。ID3算法根据信息论理论，采用划分后样本集的不确定性作为衡量划分好坏的标准，用信息增益值度量不确定性：信息增益值越大，不确定性越小。因此，ID3算法在每个非叶节点选择信息增益最大的属性作为测试属性，这样可以得到当前情况下最纯的拆分，从而得到较小的决策树。

设$S$是$s$个数据样本的集合。假定类别属性有$m$个不同的值：$C_i(i=1,2,...,m)$ 。

设$s_i$是类$C_i$中的样本数。对一个给定的样本，它总的信息熵为
\begin{equation}
I(s_1,s_2,...,s_m)=a\sum_{i=1}^m P_i-\log_2(P_i)
 \end{equation}
                  
其中,$P_i$ 是任意样本属于$C_i$的概率，一般可以用$\frac{s_i}{s}$ 估计。

设一个属性A具有$k$个不同的值，$\{a_1,a_2,...,a_k\}$利用属性A将集合$S$划分为个$k$子集，$\{S_1,S_2,...,S_k\}$其中$S_j$包含了集合$S$中属性A取$a_j$值的样本。若选择属性A为测试属性，则这些子集就是从集合$S$的节点生长出来的新的叶节点。设$s_{ij}$是子集$S_j$中类别为$C_i$的样本数，则根据属性A划分样本的信息熵值为

\begin{equation}
E(A)=a\sum_{j=1}^k\frac{s_{1j}+s_{2j}+...+s_{mj}}{s}I(s_{1j},s_{2j},...,s_{mj})
\end{equation}

其中,$I(s_{1j},s_{2j},...,s_{mj})=a\sum_{i=1}^mP_{ij}-\log_2(P_{ij})$,$P_{ij}=\frac{s_{ij}}{s_{1j}+s_{2j}+...+s_{mj}}$是子集$S_j$中类别为$C_i$的样本的概率。

最后，用属性A划分样本集S后所得的信息增益（Gain）为

\begin{equation}
 Gain(A)=I(s_1,s_2,...,s_m)-E(A) 
\end{equation}

显然E(A)越小，Gain(A)的值越大，说明选择测试属性A对于分类提供的信息越大，选择A之后对分类的不确定程度越小。属性A的k个不同的值对应的样本集S的k个子集或分支，通过递归调用上述过程（不包括己经选择的属性），生成其他属性作为节点的子节点和分支来生成整个决策树。ID3决策树算法作为一个典型的决策树学习算法，其核心是在决策树的各级节点上都用信息增益作为判断标准来进行属性的选择，使得在每个非叶节点上进行测试时，都能获得最大的类别分类增益，使分类后的数据集的熵最小。这样的处理方法使得树的平均深度较小，从而有效地提高了分类效率。

ID3算法具体流程

ID3算法的具体详细实现步骤如下：

(1)对当前样本集合，计算所有属性的信息增益；

(2)选择信息增益最大的属性作为测试属性，把测试属性取值相同的样本划为同一个子样本集；

(3)若子样本集的类别属性只含有单个属性，则分支为叶子节点，判断其属性值并标上相应的符号，然后返回调用处；否则对子样本集递归调用本算法。

下面将结合餐饮案例实现ID3的具体实施步骤。T餐饮企业作为大型连锁企业，生产的产品种类比较多，另外涉及的分店所处的位置也不同，数目比较多。对于企业的高层来讲,了解周末和非周末销量是否有大的区别，以及天气、促销活动这些因素是否能够影响门店的销量这些信息至关重要。因此，为了让决策者准确了解和销量有关的一系列影响因素，需要构建模型来分析天气、是否周末和是否有促销活动对销量的影响，下面以单个门店来进行分析。

对于天气属性，数据源中存在多种不同的值，这里将那些属性值相近的值进行类别整合。如天气为“多云”、“多云转晴”、“晴”这些属性值相近，均是适宜外出的天气，不会对产品销量有太大的影响，因此将它们为一类，天气属性值设置为“好”，同理对于“雨”、“小到中雨”等天气，均是不适宜外出的天气，因此将它们为一类，天气属性值设置为“坏”。

对于是否周末属性，周末则设置为“是”，非周末则设置为“否”。

对于是否有促销活动属性，有促销则设置为“是”，无促销则设置为“否”。

产品的销售数量为数值型，需要对属性进行离散化，将销售数据划分为“高”和“低”两类。将其平均值作为分界点，大于平均值的划分到类别“高”，小于平均值的划分为“低”类别。

经过以上的处理，我们得到的数据集合如表。

\begin{table}[thbp]\footnotesize
	\caption{处理后的数据集}
	\begin{center}
		\begin{tabular}{c|c|c|c|c}
			\hline  序号& 天气 & 是否周末&是否有促销&  销量\\ 
			\hline 1 & 坏 & 是&是&高 \\ 
			\hline 2 & 坏 & 是&是&高 \\ 
			\hline 3 & 坏 & 是&是&高 \\ 
			\hline 4 & 坏 & 否&是&高 \\ 
			\hline ...& ... & ...&...&... \\ 
			\hline 32 & 好 & 否&是&低 \\ 
			\hline 33 & 好 & 否&否&低 \\ 
			\hline 34 & 好 & 否&否&低 \\ 
			\hline
		\end{tabular} 
	\end{center}
\end{table}

采用ID3算法构建决策树模型的具体步骤如下：

(1)根据公式，计算总的信息熵，其中数据中总记录数为34，而销售数量为“高”的数据有18，“低”的有16。

\begin{equation}
I(18,16)=\frac{18}{34}\log_2\frac{18}{34}-\frac{16}{34}\log_2\frac{16}{34}=0.997503
\end{equation}

(2)根据公式，计算每个测试属性的信息熵。

对于天气属性，其属性值有“好”和“坏”两种。其中天气为“好”的条件下，销售数量为“高”的记录为11，销售数量为“低”的记录为6，可表示为(11,6)；天气为“坏”的条件下，销售数量为“高”的记录为7，销售数量为“低”的记录为10，可表示为(7,10)。则天气属性的信息熵计算过程如下：

\begin{equation}
I(11,6)=\frac{11}{17}\log_2\frac{11}{17}-\frac{6}{17}\log_2\frac{6}{17}=0.936667
\end{equation}

\begin{equation}
I(7,10)=\frac{7}{17}\log_2\frac{7}{17}-\frac{10}{17}\log_2\frac{10}{17}=0.977418
\end{equation}

\begin{equation}
E(\mbox{天气})=\frac{17}{34}I(11,6)+\frac{17}{34}I(7,10)=0.957043
\end{equation}

对于是否周末属性，其属性值有“是”和“否”两种。其中是否周末属性为“是”的条件下，销售数量为“高”的记录为11，销售数量为“低”的记录为3，可表示为(11，3)；是否周末属性为“否”的条件下，销售数量为“高”的记录为7，销售数量为“低”的记录为13，可表示为(7,13)。则节假日属性的信息熵计算过程如下：

\begin{equation}
I(11,3)=\frac{11}{14}\log_2\frac{11}{14}-\frac{3}{14}\log_2\frac{3}{14}=0.749595
\end{equation}

\begin{equation}
I(7,13)=\frac{7}{20}\log_2\frac{7}{20}-\frac{13}{20}\log_2\frac{13}{20}=0.934068
\end{equation}

\begin{equation}
E(\mbox{是否周末})=\frac{14}{34}I(11,3)+\frac{20}{34}I(7,13)=0.858109
\end{equation}

对于是否有促销属性，其属性值有“是”和“否”两种。其中是否有促销属性为“是”的条件下，销售数量为“高”的记录为15，销售数量为“低”的记录为7，可表示为（15，7）；其中是否有促销属性为“否”的条件下，销售数量为“高”的记录为3，销售数量为“低”的记录为9，可表示为(3，9)。则是否有促销属性的信息熵计算过程如下：

\begin{equation}
I(15,7)=\frac{15}{22}\log_2\frac{15}{22}-\frac{7}{22}\log_2\frac{7}{22}=0.902393
\end{equation}

\begin{equation}
I(3,9)=\frac{3}{12}\log_2\frac{3}{12}-\frac{9}{12}\log_2\frac{9}{12}=0.811278
\end{equation}

\begin{equation}
E(\mbox{是否有促销})=\frac{22}{34}I(15,7)+\frac{12}{34}I(3,9)=0.870235
\end{equation}

(3)根据公式，计算天气、是否周末和是否有促销属性的信息增益值。

$Gain(\mbox{天气})=I(18,16)-E(\mbox{天气})=0.997503-0.957043=0.04046$

$Gain(\mbox{是否周末})=I(18,16)-E(\mbox{是否周末})=0.997503-0.858109=0.139394$

$Gain(\mbox{是否有促销})=I(18,16)-E(\mbox{是否有促销})=0.997503-0.870235=0.127268$

(4)由第3步的计算结果可以知道是否周末属性的信息增益值最大，它的两个属性值“是”和“否”作为该根结点的两个分支。然后按照第1步到第3步所示步骤继续对该根结点的三个分支进行结点的划分，针对每一个分支结点继续进行信息增益的计算，如此循环反复，直到没有新的结点分支，最终构成一棵决策树。生成的决策树模型如下图：

\begin{figure}[thbp!]
	\centering
	\includegraphics[width=0.4\linewidth]{figure/3-1}
	\caption{\ 全局离群点和局部离群点}
	\label{fig:3-1}
\end{figure}

从上面的决策树模型可以看出，门店的销售高低和各个属性之间的关系，并可以提取出以下决策规则：

\begin{itemize}
	\item 若周末属性为“是”，天气为“好”，则销售数量“高”；
	\item 若周末属性为“是”，天气为“坏”，促销属性为“是”，则销售数量“高”；
	\item 若周末属性为“是”，天气为“坏”，促销属性为“否”，则销售数量“低”；
	\item 若周末属性为“否”，促销属性为“否”，则销售数量“低”；
	\item 若周末属性为“否”，促销属性为“是”，天气为“好”，则销售数量“高”；
	\item 若周末属性为“否”，促销属性为“是”，天气为“坏”，则销售数量“低”；
\end{itemize}

由于ID3决策树算法采用了信息增益作为选择测试属性的标准，会偏向于选择取值较多的即所谓高度分支属性，而这类属性并不一定是最优的属性。同时ID3决策树算法只能处理离散属性，对于连续型的属性，在分类前需要对其进行离散化。为了解决倾向于选择高度分支属性的问题，人们采用信息增益率作为选择测试属性的标准，这样便得到C4.5决策树算法。此外常用的决策树算法还有CART算法、SLIQ算法、SPRINT算法和PUBLIC算法等等。

\subsection{数据决策树处理}

本文主要使用RapidMiner中的决策树算子对所给数据进行处理后分别生产两个决策树：

(1)驻地决策树：这个决策树是以文化程度为label，以驻地为展示属性来反映不同文化程度数据类构造的驻地规则。

(2)工种决策树：这个决策树是以工种为label，反映出不同工种对应的文化程度，驻地的相关性规律。

	\section{已知问题和未来发展}
\subsection{已知问题}
本模板未采用2013版规范的页边距设置，因为实在是办不到2.5CM顶部页边距加上2.6CM的页眉设置啊。
\subsection{未来发展}
西安石油大学本科生论文的未来发展还是需要各位用户的参与，如果每一个用户都能贡献出一点关于\LaTeX 模板的想法和意见，我相信几年之后西安石油大学本科生论文模板会成为其他高校学习和借鉴的例子。同学当自强，让我们一起来丰富完善这个模板，如果你有很好的建议或者意见请发送到rikrun@gmail.com
\subsection{官方认证}
到目前为止（\today ）没有西安石油大学任何官方组织对于本模板的格式或者内容进行认证，这代表采用本模板进行的论文写作可能不被官方的论文系统接受。
	\section{系统开发成果}

对西安石油大学计算机学院 GMS 系统进行的第二期工程已完成设计与开发，通过测试并于 2018 年 5 月运行在西安石油大学内，校园网网址：\underline{http://bkbysj.xsyu.edu.cn/}。将承接计算机学院 2018 年毕业生毕业设计相关线上工作。

\subsection{毕业设计答辩模块}

整个毕业设计答辩阶段的流程从管理员划分答辩小组开始，该划分操作逻辑复杂，既要考虑组内教师的职称、工作年限，又要综合考虑教师今年所指导的学生（尽量使学生的指导教师不在其答辩组内）。所以，系统将划分答辩小组的操作交由管理员进行，线上系统则提供辅助决策的作用。

管理员在如图 \ref{divide-show} 所示的页面中进行答辩小组划分。首先下载一个由系统生成的 Excel 文档，该文档中包含由答辩系统预生成的小组划分表。然后以该表作为参考，管理员对小组划分进行微调，完成后将该表回传至系统内，即可完成当年的答辩小组划分。

\begin{figure}
	\centering
	\includegraphics[width=1\linewidth]{figure/divide-show}
	\caption{管理员答辩小组划分页面}
	\label{divide-show}
\end{figure}

在管理员完成答辩小组划分之后的下一个流程，是指导教师对自己所带学生进行论文审阅。进入图 \ref{tealook} 所示页面后，依次对学生们点击“进行论文审阅”按钮，输入分数和内容即可完成审阅操作，并下载论文审阅书。这时该学生的“进行论文审阅”按钮会变成“√修改论文审阅”，再次点击该按钮可对已审阅的内容进行修改。

该页面同时是每位教师对自己所带学生的管理页面（关键信息已模糊化处理）。值得一说的是，本系统中所有的用户头像均不相同，根据用户 UID 哈希生成的随机矢量风格，避免了所有老师学生都使用默认头像的尴尬又无聊的景象。

\begin{figure}
	\centering
	\includegraphics[width=\linewidth]{figure/5-1}
	\caption{教师管理所带学生页面}
	\label{tealook}
\end{figure}

审阅阶段过后，答辩流程将会进入到评阅环节，教师进行答辩评阅的页面如图 \ref{tea-review} 所示。添加评阅时，教师在线下阅读学生的论文及相关资料，然后在系统中输入该学生的学号、评阅内容、评阅成绩，即可完成评阅操作。为了方便教师操作，可支持仅输入学号后数位，系统会自动匹配到对应学生并完成评阅操作，如果存在相同学号后缀的学生，系统亦会给与提示。

从截图中可以看出，系统同时显示出了教师所属的班级，以及小组的答辩地点，这是在深度了解教师需求后开发的人性化功能。评阅教师在完成所有评阅数据的录入后，可点击“生成评阅意见表”按钮，会下载得到一个压缩包，内包含评阅的所有学生各自的评阅意见表，该表符合学校相关的格式要求，可直接提供给学校教务处。

\begin{figure}
	\centering
	\includegraphics[width=\linewidth]{figure/tea-review}
	\caption{答辩组成员评阅页面}
	\label{tea-review}
\end{figure}

在此系统中小组秘书拥有最多的线上功能，包括统计组内成员工作量以供生成工作量统计表，以及小组内学生的答辩成绩、总成绩的录入。秘书的成绩录入工作使用 Excel 辅助进行，而非在网页上逐次添加。秘书首先下载一个系统自动生成的小组内学生名单，其中包括学生的审阅成绩和评阅成绩。只需要输入答辩成绩之后，总成绩会自动按照公式生成。最后再将填写完成后的 Excel 文档上传，即可完成组内学生成绩的录入。

完成答辩阶段所有数据的录入之后，管理员和答辩秘书、答辩组长同时拥有类似图 \ref{data-show} 的数据展示页面，以表格的形式展现出对应权限所拥有的学生、题目等数据和信息，表格也可以很方便的自定义按列排序、搜索关键字，以及导出为 Excel、PDF 等格式。

\begin{figure}
	\centering
	\includegraphics[width=\linewidth]{figure/5-3}
	\caption{答辩相关数据展示页面}
	\label{data-show}
\end{figure}


\subsection{系统内其他辅助模块}

专门为管理员设计的功能占整个系统工作量的 60\% 以上。在管理员面板中，可以管理所有的用户，设定每位用户的类型，也可以看到现在所有选题配对情况。值得一说的是，管理员支持使用 Excel 批量导入每届学生老师信息，系统会自动解析 xlsx 文档，并创建对应的登陆账号。此外，本系统网页中所有可见的表格信息均可一键导出为 Excel 或 Word 文档，方便进一步办公处理。

图 \ref{admin-panel} 是管理员对全院学生综合管理的仪表盘页面，在这里利用了数据可视化技术，将全院所有学生的综合情况清晰地展现在管理员面前。在这里，可以直观的看到待选题目和学生总数的柱状对比图，也可以分专业以饼状图的形式看到当前各专业学生的选题状态分布。这些都是选题流程中管理员需要掌握的数据。

\begin{figure}
	\centering
	\includegraphics[width=0.85\linewidth]{figure/5-2}
	\caption{管理员仪表盘页面}
	\label{admin-panel}
\end{figure}

学生作为本系统中最简单的角色，可进行选题，以及在选题成功之后通过此系统向老师发送文件。在选题方面设计了 2 个人性化的特性：能看到某道题当前已选人数，这大大避免某道题被大家集中选择；在教师查看你的选题志愿之前，可以取消申请，并另选一道题。（事实上在此系统中所有的状态转移均支持最大程度的撤销操作）

该系统中的教师角色可以出题，并且实时跟踪自己题目的状态，历年所出题目会形成一个自己的题库以供复用，题目支持上传附件。这些特性弥补了计算机学院旧选题系统的遗憾。

本系统还拥有一个强大的自动备份还原功能，系统会自动在每天凌晨 3 时进行一次数据库备份，同时自动删除 15 天前的备份（不支持手动删除），当然，用户可以选择在需要的时候随时手动创建一个备份。这样的设计使得系统更加稳定，无论是管理员的误操作，还是被任何形式的恶意攻击，都不会对系统造成很大的影响。

此外，如图 \ref{gms-news} 所示，本系统拥有完善的公告系统，支持富文本编辑、设置置顶、支持设置公告对不同类型用户的可见性，以及附件支持。

\begin{figure}
	\centering
	\includegraphics[width=1\linewidth]{figure/gms-news}
	\caption{GMS 公告系统}
	\label{gms-news}
\end{figure}

	\section{系统测试}

软件或硬件的系统测试是在完整的集成系统上进行测试，以评估系统是否符合其指定的要求。通常，系统测试将所有已通过集成测试的软件组件作为其输入，并将软件系统本身与适用的硬件系统集成在一起。集成测试的目的是检测集成在一起的软件单元（称为集合）和硬件之间的任何不一致。系统测试是一种有限的测试类型，它试图检测整个系统内的缺陷。

\subsection{黑盒测试}

黑盒测试是一种软件测试方法，它可以检查应用程序的功能，而无需查看其内部结构或工作情况。这种测试方法可以直接应用在各个级别的软件测试当中：单元测试、集成测试、系统测试和系统验收。所以它有时被称为基于规格的测试。

一般情况下，是不需要应用程序代码/内部结构和编程知识的具体知识的。测试人员知道该软件应该做什么，但不知道它是如何做到的。例如，测试人员知道特定的输入会返回一定的不变的输出，但不知道软件如何产生输出。

为了实现黑盒测试，针对答辩系统编写了测试用例。测试用例是围绕规范和要求构建的，即应用程序应该做什么。测试用例一般都来自于软件的外部描述，包括设计规范、需求要求和设计参数等。虽然使用的测试主要是功能性的，但也可以使用非功能测试。测试设计人员选择有效和无效的输入，并在测试人员的帮助下或者已知结果良好的情况下确定正确的输出，而无需了解测试对象的内部结构。

\subsection{负载压力测试}

负载测试在开发社区中以不同的方式被使用。负载测试通常是指通过模拟多个用户同时访问程序来对软件程序的预期使用效果进行预估。因此，这种测试与多用户系统最相关，例如针对该答辩项目的Web服务器进行的压力测试。最准确的负载测试为模拟实际使用，而不是使用理论或分析建模进行测试。

通过负载测试，我们可以根据实际的客户行为来衡量网站的服务质量（QOS）性能。几乎所有的负载测试工具和框架都遵循经典的负载测试范例：当客户访问您的网站时，脚本记录器记录通信，然后创建相关的交互脚本。负载生成器尝试重放脚本，这些脚本可能在重放时使用不同的测试参数进行调用。在重放过程中，硬件和软件统计数据将被监视和收集，这些统计数据包括物理服务器的CPU、内存、磁盘IO以及响应时间、被测系统（SUT）的吞吐量等，最后综合分析所有这些统计数据，并生成负载测试报告。

在这里，我们不能使用传统的测试工具，比如 ApacheBench（一个单线程命令行计算机程序，用于测量HTTP Web服务器的性能。最初设计用于测试Apache HTTP Server）。因为Meteor代码不会在每个请求上执行，而是首先在浏览器端运行。如果我们发送连续的压力测试请求，将会得到静态的JS代码，这对评估没有意义。

于是在发布到生产环境之前，我们使用MeteorDown为我们的答辩系统进行压力测试。它使用DDP协议与Meteor应用程序进行通信。首先，150行数据被添加到MongoDB的title集合中，每个title文档包含14个字段。在每次测试中，MeteorDown不仅会向用户请求，还会通过Meteor核心提供的方法订阅150行数据。所需的数据量非常大，超出了正常使用的范围。足以反映系统在极高负载下的稳定性。

测试环境的配置信息如表6-1所示，最终的测试结果如图6-1所示。通过测试结果我们可以看到，随着并发数量的增加，虽然单次响应时间逐渐提升，但始终拥令人满意的较高响应速度。每次测试需要与服务器端  Meteor 建立连接，请求并下载数据库中 150 行复杂的 documents。Meteor 可以对这样的测试每分钟运行近万次，并且随着并发数的提高，性能并没有出现明显的下降。这样的性能足以部署在生产环境中，满足大量校园用户在同一时间使用该系统的需求。

\begin{table}
	\centering
	\caption{测试环境配置}  %表格标题
	\begin{tabular}{ll} 
		\hline
		\hline
		核心硬件 & 型号 \\ 
		\hline
		CPU	& Intel Core i7-4700MQ CPU（2.40GHz）\\
		RAM	& DDR3 1600MHz 4Gb*2 \\
		SSD	& PLEXTOR SATA3 PX-128M6M \\
		\hline
		\hline
		核心软件 & 版本 \\ 
		\hline
		OS	& Windows 10 64-bit Professional Edition \\
		MeteorDown & v2.6.0 \\
		Node.js & v4.6.1 \\
		Meteor & v1.4.2 \\
		\hline
	\end{tabular}
\end{table}

\begin{figure}
	\centering
	\includegraphics[width=0.7\linewidth]{figure/test_result}
	\caption{压力测试结果}
\end{figure}

负载测试的表现明显超过使用传统后端程序生成 HTML 页面的系统，比如 PHP、ASP.NET 以及 JSP 等编程语言。这这说明服务器端只传输数据，而将数据的渲染交给客户端，能有效降低服务器的运行负载。

	\section*{结论}
\addcontentsline{toc}{section}{结论}

根据开发高校毕业设计答辩管理系统的需求，我们首先将功能抽象化，确定了毕业设计答辩工作的流程，以及用户角色划分。采用了JS全栈开发和Meteor开发框架，数据库采用 MongoDB，以及将React作为前端数据渲染引擎。这些都是目前Web开发中最为流行的工具。Node.js作为一个将JavaScript从浏览器带到服务器的项目，现在已经获得肯定并广泛应用于Web开发中。而问答社区Stack Overflow 2016年开发者调查报告中显示，JavaScript 再一次成为了全球最受欢迎的开发技术。

除了可以直接使用来自众多开源工作者的开源项目，JS应用的开发过程也逐渐工业化。包括使用CSS的预处理语言LESS，通过 Babel 工具也能将使用了最新规范ES2015的代码转换为浏览器普遍兼容的ES5代码。将CSS和JS代码批量合并、混淆、压缩，获得了网页更快的加载速度。使用ESLint能将代码出错的可能性降低到最小，并且拥有统一的编码风格，有利于团队协作。同时也介绍了被使用在我们项目中的基于 Git 的工作流。

我们遵循了Meteor和React所提倡的理念：前后端之间只对数据进行交互，对数据的渲染由客户端完成而非服务器端。包括对数据库读写的代码，都在客户端中被执行，节省了发送和接受请求的过程，同时安全性也是有保证的。在实际的开发中，我们感受到了这些特性所带来的便利性，那就是应用逻辑清晰，并且开发速度很快。也正是由于将数据的渲染交给了浏览器，从而大大降低了服务器的负载，在最后的压力测试中表现出了很好的成绩。

	
	%============= 参考文献 =====================
	\addcontentsline{toc}{section}{参考文献}
	\nocite{*}
	\bibliographystyle{gbt7714-2005}  % 论文引用格式
	\zihao{5} 
	\bibliography{bibfile}
	\zihao{-4} 
	
	%=============  致谢  ======================
	\section*{致谢}
\addcontentsline{toc}{section}{致谢}
本次论文是在我的导师张留美张老师的亲切关心和悉心指导下完成的。他指导学生细心、认真，大到整个设计思路是否正确，小到一个数据属性字段的含义。从论文课题的选择到最终完成，张老师都始终给予我细心的指导和不懈的支持。在此谨向张老师致以诚挚的谢意和崇高的敬意。

此外，我还要感谢我的舍友，是他们开朗的性格，以及活跃的解决问题思路，让我找到算法设计的初模。当然，他们也在闲暇之余，和我一起讨论一些算法的思路，让我受益匪浅。在这，对你们说声谢谢，这份舍友情将一直会陪伴着我。

在我的十几年求学历程里，离不开父母的鼓励和支持，是他们辛勤的劳作，无私的付出，为我创造良好的学习条件，我才能顺利完成完成学业，感激他们一直以来对我的抚养与培育。

或许，毕业设计，将是我大学生涯交上的最后一个作业了。真的想借此机会对四年以来给我帮助的所有老师、同学、朋友说声谢谢，你们的友谊将是我人生的最大的财富，是我生命中不可或缺的一部分。

在整个系统的设计过程中，遇到过很多障碍，挫折，然而当整个系统完整的展现在自己的面前时，这种喜悦是只可意会的。几个月来忙碌紧张而又有条不紊的毕业设计，使我有机会对本专业的基本理论、专业知识和基本技术有了更深入的了解和体会，使我在大学中所学到的知识得到了系统和升华，真正达到了学以致用。

这次做论文的经历会使我终身受益，我感受到做一件事情，从开始到结束的那份收获。也让我明白认真做一件事情的重要性。希望这次的经历能让我在以后学习中激励我继续进步。

	
	%=============  附录 ======================
	%\section*{附录}
\addcontentsline{toc}{section}{附录}
在本次论文的写的过程，会用到以下算法，现我将简要将这些算法介绍。

\title{附录1  算法类似性}

(1)决策树学习：根据数据的属性采用树状结构建立决策模型。决策树模型常常用来解决分类和回归问题。常见的算法包括 CART (Classification And Regression Tree)、ID3、C4.5、随机森林 (Random Forest) 等。

(2)回归算法：试图采用对误差的衡量来探索变量之间的关系的一类算法。常见的回归算法包括最小二乘法 (Least Square)、逻辑回归 (Logistic Regression)、逐步式回归 (Stepwise Regression) 等。

(3)聚类算法：通常按照中心点或者分层的方式对输入数据进行归并。所有的聚类算法都试图找到数据的内在结构，以便按照最大的共同点将数据进行归类。常见的聚类算法包括 K-Means 算法以及期望最大化算法 (Expectation Maximization) 等。

(4)人工神经网络：模拟生物神经网络，是一类模式匹配算法。通常用于解决分类和回归问题。人工神经网络算法包括感知器神经网络 (Perceptron Neural Network) 、反向传递 (Back Propagation) 和深度学习等。

\newpage
\title{附录2   决策树模型评估}

建立了决策树模型后需要给出该模型的评估值，这样才可以来判断模型的优劣。学习算法模型使用训练集 (training set) 建立模型，使用校验集 (test set) 来评估模型。本文通过评估指标和评估方法来评估决策树模型。 评估指标有分类准确度、召回率、虚警率和精确度等。而这些指标都是基于混淆矩阵 (confusion matrix) 进行计算的。

混淆矩阵是用来评价监督式学习模型的精确性，矩阵的每一列代表一个类的实例预测，而每一行表示一个实际的类的实例。以二类分类问题为例，如下表所示：

\begin{table}[thbp]\footnotesize
	\caption{混淆矩阵}
	\begin{center}
		\begin{tabular}{cc|clc}
			\hline	&预测的类\\
			  实际的类& & 类=1 & 类=0  \\ 
			\hline &类=1 & TP & FN &P\\ 	
			\hline &类=0 & FP & TN&N \\ 
			\hline
		\end{tabular} 
	\end{center}
\end{table}


P (Positive Sample)：正例的样本数量。

N(Negative Sample)：负例的样本数量。

TP(True Positive)：正确预测到的正例的数量。

FP(False Positive)：把负例预测成正例的数量。

FN(False Negative)：把正例预测成负例的数量。

TN(True Negative)：正确预测到的负例的数量。

根据混淆矩阵可以得到评价分类模型的指标有以下几种。

分类准确度，就是正负样本分别被正确分类的概率，计算公式为：

\begin{equation}
Accuracy=\frac{TP+TN}{P+N}
\end{equation}

召回率，就是正样本被识别出的概率，计算公式为：

\begin{equation}
Recall=\frac{TP}{P}
\end{equation}

虚警率，就是负样本被错误分为正样本的概率，计算公式为：

\begin{equation}
FPrate=\frac{FP}{N}
\end{equation}

精确度，就是分类结果为正样本的情况真实性程度，计算公式为：

\begin{equation}
Precision=\frac{TP}{TP+FP}
\end{equation}

评估方法有保留法、随机二次抽样、交叉验证和自助法等。

保留法 (holdout) 是评估分类模型性能的最基本的一种方法。将被标记的原始数据集分成训练集和检验集两份，训练集用于训练分类模型，检验集用于评估分类模型性能。但此方法不适用样本较小的情况，模型可能高度依赖训练集和检验集的构成。

随机二次抽样 (random subsampling) 是指多次重复使用保留方法来改进分类器评估方法。同样此方法也不适用训练集数量不足的情况，而且也可能造成有些数据未被用于训练集。

交叉验证 (cross-validation) 是指把数据分成数量相同的 k 份，每次使用数据进行分类时，选择其中一份作为检验集，剩下的 k-1 份为训练集，重复 k 次，正好使得每一份数据都被用于一次检验集 k-1 次训练集。该方法的优点是尽可能多的数据作为训练集数据，每一次训练集数据和检验集数据都是相互独立的，并且完全覆盖了整个数据集。也存在一个缺点，就是分类模型运行了 K 次，计算开销较大。

自助法 (bootstrap) 是指在其方法中，训练集数据采用的是有放回的抽样，即已经选取为训练集的数据又被放回原来的数据集中，使得该数据有机会能被再一次抽取。用于样本数不多的情况下，效果很好。
\appendix
	
\end{document}
%%%%%%%%%% 结束 %%%%%%%%%%
